\begin{claim}{C.88}{Proved}{The exact vanishing-displacement image and its compact-family topology}
Retain all the actual hypotheses and maps of
C.81: the original closed disk $\overline D$,
its interior $O$, a compatible metric $d$
on $\overline D$, a closed discrete free
proper operator group $Q$ on $O$, and a
transitive subgroup of the original
boundary-fixed centralizer $\mathscr C_Q$.
Write $\pi:O\to N=O/Q$ and retain
$d_Q(x,y)=\sup_{q\in Q}d(qx,qy)$.
For an interior homeomorphism $h$ commuting
with every actual $q$, put
\[
 \delta_h(\pi x)=d_Q(hx,x),\qquad
 \mathscr V_Q=
 \{h\in C_{\operatorname{Homeo}(O)}(Q):
                         \delta_h\in C_0(N)\}.                 \tag{C.88.1}
\]
Here $C_0(N)$ means continuous real-valued
functions whose absolute values are uniformly
small off compact subsets of $N$.
Then restriction and identity extension are
mutually inverse topological group maps
\[
 \begin{split}
 R:\mathscr C_Q&\longrightarrow\mathscr V_Q,
                       &R(a)&=a|_O,\\
 E:\mathscr V_Q&\longrightarrow\mathscr C_Q,
                       &E(h)|_O&=h,\quad
                         E(h)|_{\partial\overline D}=\id .
 \end{split}                                                   \tag{C.88.2}
\]
The topology on $\mathscr V_Q$ is the sum
of the uniform $d_Q$ distances of maps and
inverses. These maps are isometries for the
corresponding original closed-disk distances.
Every member of $\mathscr V_Q$ and its inverse
is automatically uniformly $d_Q$-continuous;
no uniform-continuity assumption is added
to (C.88.1).

There is also an exact compact-parameter
criterion. Let $P$ be a compact metric space
and let $h_p$ be actual equivariant interior
homeomorphisms such that $(p,x)\mapsto h_p(x)$
is continuous on $P\times O$. Their identity
extensions form a continuous $P$-family of
original closed-disk homeomorphisms and
inverses precisely when
\[
 \forall\epsilon>0\ \exists K\subset N
 \text{ compact}\quad
 \sup_{\substack{p\in P\\ \pi x\notin K}}
                    d_Q(h_p x,x)<\epsilon.                    \tag{C.88.3}
\]
In that case the inverse family has the
same uniform vanishing property. Thus
(C.88.3), with the actual interior maps
retained, expresses the original boundary
and compact-family requirements exactly.
It supplies no cutoff deformation by itself.
\end{claim}

\begin{proof}
\emph{The actual displacement functions and algebra.}
C.81 proves that $d_Q$ induces the original
topology of $O$, that $\pi$ is the specified
covering onto the locally compact surface
$N$, and that $d_Q$ is $Q$-invariant.
Consequently $x\mapsto d_Q(hx,x)$ is
continuous. Commutation and invariance give
\[
 d_Q(hqx,qx)=d_Q(qhx,qx)=d_Q(hx,x).
                                                                  \tag{C.88.4}
\]
The quotient map therefore gives exactly
the continuous function in (C.88.1).
The map $h$ descends to the homeomorphism
$\bar h(\pi x)=\pi(hx)$, whose displayed
inverse is $\overline{h^{-1}}$.
Symmetry and the triangle inequality give
\[
 \begin{split}
 \delta_{h^{-1}}&=\delta_h\circ\bar h^{-1},\\
 \delta_{hk}&\leq\delta_h\circ\bar k+\delta_k .
 \end{split}                                                   \tag{C.88.5}
\]
A homeomorphism of $N$ and its inverse
carry compact sets to compact sets.
Thus its pullback preserves $C_0(N)$.
Equations (C.88.5) prove that
$\mathscr V_Q$ is a group, before using
the boundary extension.

\emph{Necessity, including the common tail for compact families.}
Let $a_p$, $p\in P$, be a continuous compact
family in $\mathscr C_Q$. Joint evaluation
on $P\times\overline D$ is continuous.
Since every $a_p$ fixes every boundary
point, for any $\epsilon>0$ there is a
compact set $C\subset O$ such that
\[
          d(a_p z,z)<\epsilon
                \quad(p\in P,\ z\in O\setminus C).             \tag{C.88.6}
\]
For completeness, failure would give
$p_j\to p$ after taking a subsequence
and $z_j$ escaping all interior compact
sets. A further subsequence converges
to $z\in\partial\overline D$, and joint
continuity gives $a_{p_j}z_j\to z$,
contradicting the displayed failure.

The set $K=\pi C$ is compact. If
$\pi x\notin K$, then $qx\notin C$
for every $q\in Q$. The original
commutator equations imply
\[
 \begin{split}
 d_Q(a_p x,x)
   &=\sup_{q\in Q}d(a_p qx,qx)\\
   &\leq\epsilon
       \qquad(p\in P,\ \pi x\notin K).
 \end{split}                                                   \tag{C.88.7}
\]
Starting with $\epsilon/2$ gives the
strict bound in (C.88.3).
The same proof applies to the compact
inverse family. Taking $P$ to be a
point proves that $R$ maps into
$\mathscr V_Q$.

\emph{One common control for compact quotient cores.}
We recall the exact C.81 comparison
needed for the reverse direction.
Let $C,C'\subset O$ be compact.
For a fixed $x_0\in O$, openness of
evaluation for the actual transitive
Polish source, and its compact-covering
lemma, give compact families $S,T$
of original boundary-fixed maps with
$C\subset S(x_0)$ and $C'\subset T(x_0)$.
For $c=sx_0\in C$ and $c'=tx_0\in C'$,
commutation gives
\[
          qc'=tqx_0=ts^{-1}(qc).
                                                                  \tag{C.88.8}
\]
All choices on the right lie in the
single compact actual family $TS^{-1}$.
It follows that, for $c_j\in C$ and
$q_jc_j\to z\in\partial\overline D$,
\[
                   q_j C'\longrightarrow\{z\}
                  \quad\hbox{uniformly}.                       \tag{C.88.9}
\]
Indeed joint evaluation of the compact
family $TS^{-1}$ fixes $z$ throughout.
No limiting map or original-circle
extension of the operators $q_j$
has been used.

\emph{Sufficiency and all parameter and inverse maps.}
Suppose the interior family in the
statement satisfies (C.88.3), and define
$E(h_p)$ by the literal formula (C.88.2).
Interior continuity is given. To prove
joint continuity at the original circle,
take $p_j\to p$ and $x_j\to z$ with
$z\in\partial\overline D$.
Terms already on the circle are fixed.
If the desired convergence fails for
the remaining interior terms, retain
a sequence on which it fails by a
fixed positive distance.

The surface $N$ has a compact exhaustion.
After passage to a subsequence, either
$\pi x_j$ leaves every compact subset
of $N$, or all $\pi x_j$ belong to
one compact $K\subset N$. This
dichotomy follows by taking infinitely
many terms in one exhaustion set when
there is such a set; otherwise the
sequence itself leaves every set.
In the escaping case, (C.88.3) gives
\[
             d(h_{p_j}x_j,x_j)
                  \leq d_Q(h_{p_j}x_j,x_j)\longrightarrow0,
                                                                  \tag{C.88.10}
\]
which contradicts failure.

In the compact-core case choose a
compact $C\subset O$ with $\pi C=K$.
Finitely many local covering sections
over relatively compact neighborhoods
of points of $K$ give such a $C$.
Write $x_j=q_jc_j$, $c_j\in C$.
The joint interior continuity and
compactness of $P\times C$ make
\[
                 C'=\{h_p(c):p\in P,\ c\in C\}
                                                                  \tag{C.88.11}
\]
compact in $O$. By (C.88.9),
$q_j C'\to\{z\}$ uniformly.
But $h_{p_j}x_j=q_jh_{p_j}c_j$
belongs to $q_j C'$, again a
contradiction. This proves joint
boundary continuity in every case.

Each $E(h_p)$ is bijective: its
interior restriction is a homeomorphism
onto $O$ and its boundary restriction
is the identity. The actual map
\[
 \mathcal E:P\times\overline D
       \longrightarrow P\times\overline D,\qquad
             (p,x)\longmapsto(p,E(h_p)x)                         \tag{C.88.12}
\]
is consequently a continuous bijection
between compact Hausdorff spaces.
Its inverse is continuous and has
the exact formula
$(p,y)\mapsto(p,E(h_p)^{-1}y)$.
Thus both the family and its inverses
are continuous in the original
closed-disk topology. The original
operator equations hold on $O$
unchanged. Applying the proved
necessity to this inverse family
gives its uniform $C_0$ tails.
This proof even obtains joint inverse
continuity without assuming it for
the initial interior family.

With $P$ a point, this establishes
that $E$ maps every member of
$\mathscr V_Q$ into $\mathscr C_Q$.
The two formulas (C.88.2) compose
to the identity on their exact
domains. Extension commutes with
composition and inversion, since
these identities hold on $O$ and
both sides fix the original circle.

\emph{The complete topology comparison.}
A modulus of uniform continuity
of $E(h)$ on $\overline D$ applies
to each pair $qx,qy$. Commutation
then proves uniform continuity of
$h$ for $d_Q$. Apply the same
argument to $E(h)^{-1}$.
For $h,k\in\mathscr V_Q$, the
actual equality is
\[
 \begin{split}
 \sup_{x\in O}d_Q(hx,kx)
  &=\sup_{\substack{x\in O\\ q\in Q}}
               d(E(h)qx,E(k)qx)\\
  &=\sup_{z\in\overline D}d(E(h)z,E(k)z).
 \end{split}                                                   \tag{C.88.13}
\]
The last equality uses $q=1$ and
density of $O$. Repeating it for
both inverses proves the asserted
isometry and both continuous maps.
In particular the displacement map
to $C_0(N)$ satisfies
\[
 \|\delta_h-\delta_k\|_\infty
       \leq\sup_{x\in O}d_Q(hx,kx).                             \tag{C.88.14}
\]
All image and topology requirements
are therefore expressed using the
specified covering and the original
operator labels.

The remaining general cutoff problem
is to deform an arbitrary original
based sphere family to a family with
one compact quotient support while
retaining joint interior continuity,
all equivariance equations and the
uniform tail condition (C.88.3).
C.81 fills the resulting supported
family. C.88 provides the exact
boundary criterion for such a
deformation; it does not construct
that deformation or replace its
uniform tail control by convergence
on compact subsets of $N$.
\end{proof}
